
\documentclass{article}
\usepackage{colortbl}
\usepackage{makecell}
\usepackage{multirow}
\usepackage{supertabular}

\begin{document}

\newcounter{utterance}

\centering \large Interaction Transcript for game `hot\_air\_balloon', experiment `air\_balloon\_survival\_en\_reasoning off\_easy', episode 0 with Qwen3{-}32B{-}t0.0.
\vspace{24pt}

{ \footnotesize  \setcounter{utterance}{1}
\setlength{\tabcolsep}{0pt}
\begin{supertabular}{c@{$\;$}|p{.15\linewidth}@{}p{.15\linewidth}p{.15\linewidth}p{.15\linewidth}p{.15\linewidth}p{.15\linewidth}}
    \# & $\;$A & \multicolumn{4}{c}{Game Master} & $\;\:$B\\
    \hline

    \theutterance \stepcounter{utterance}  
    & & \multicolumn{4}{p{0.6\linewidth}}{
        \cellcolor[rgb]{0.9,0.9,0.9}{
            \makecell[{{p{\linewidth}}}]{
                \texttt{\tiny{[P1$\langle$GM]}}
                \texttt{You are participating in a collaborative negotiation game.} \\
\\ 
\texttt{Together with another participant, you must agree on a single set of items that will be kept. Each of you has your own view of how much each item matters to you (importance). You do not know how the other participant values the items. Additionally, you are given the effort each item demands.} \\
\texttt{You may only agree on a set if the total effort of the selected items does not exceed a shared limit:} \\
\\ 
\texttt{LIMIT = 4379} \\
\\ 
\texttt{Here are the individual item effort values:} \\
\\ 
\texttt{Item effort = \{"A19": 862, "B91": 252, "A30": 515, "B11": 571, "C47": 211, "C75": 339, "C26": 946, "C73": 347, "A47": 522, "A15": 803, "A75": 401, "A39": 980, "A56": 918, "C24": 599, "B26": 493\}} \\
\\ 
\texttt{Here is your personal view on the importance of each item:} \\
\\ 
\texttt{Item importance values = \{"A19": 138, "B91": 583, "A30": 868, "B11": 822, "C47": 783, "C75": 65, "C26": 262, "C73": 121, "A47": 508, "A15": 780, "A75": 461, "A39": 484, "A56": 668, "C24": 389, "B26": 808\}} \\
\\ 
\texttt{Goal:} \\
\\ 
\texttt{Your goal is to negotiate a shared set of items that benefits you as much as possible (i.e., maximizes total importance to YOU), while staying within the LIMIT. You are not required to make a PROPOSAL in every message {-} you can simply negotiate as well. All tactics are allowed!} \\
\\ 
\texttt{Interaction Protocol:} \\
\\ 
\texttt{You may only use the following structured formats in a message:} \\
\\ 
\texttt{PROPOSAL: \{'A', 'B', 'C', …\}} \\
\texttt{Propose keeping exactly those items.} \\
\texttt{REFUSE: \{'A', 'B', 'C', …\}} \\
\texttt{Explicitly reject opponent's proposal.} \\
\texttt{ARGUMENT: \{'...'\}} \\
\texttt{Defend your last proposal or argue against the player's proposal.} \\
\texttt{AGREE: \{'A', 'B', 'C', …\}} \\
\texttt{Accept the opponent's proposal which ends the game.} \\
\\ 
\\ 
\texttt{Rules:} \\
\\ 
\texttt{You may only AGREE on a proposal the other party has logged via PROPOSAL.} \\
\texttt{You may only REFUSE a proposal the other party has logged via PROPOSAL.} \\
\texttt{Total effort of any PROPOSAL or AGREE set must be ≤ LIMIT.} \\
\texttt{Do NOT reveal your hidden importance scores.} \\
\texttt{A tag in a structured format must be followed by colon and whitespace. The argument must be a python set containing 0 or more strings.} \\
\texttt{So, it must be of the form TAG: \{...\}} \\
\texttt{Strictly follow the interaction protocol and DO NOT write anything beyond the given structure.} \\
\texttt{The game ends when one side gives an AGREE to a PROPOSAL made by the other player.} \\
\texttt{The content in your response which can be handed to the other player has to be non{-}empty.} \\
\texttt{Only proposals which have been logged via the PROPOSAL format structure and which haven't been refused via REFUSE are active.} \\
\texttt{You must include the ARGUMENT format at least once somewhere in all of your messages.} \\
            }
        }
    }
    & \\ \\

    \theutterance \stepcounter{utterance}  
    & \multicolumn{4}{p{0.6\linewidth}}{
        \cellcolor[rgb]{0.9,0.9,0.9}{
            \makecell[{{p{\linewidth}}}]{
                \texttt{\tiny{[P1$\rangle$GM]}}
                \texttt{<think>} \\
\texttt{Okay, let's see. I need to negotiate a set of items that maximizes my total importance without exceeding the effort limit of 4379. First, I should figure out which items are most important to me and check their combined effort.} \\
\\ 
\texttt{Looking at my importance values: B91 is 583, A30 is 868, B11 is 822, C47 is 783, A15 is 780, A75 is 461, B26 is 808, A19 is 138, A47 is 508, A39 is 484, A56 is 668, C26 is 262, C75 is 65, C73 is 121, C24 is 389. } \\
\\ 
\texttt{The highest importance items are A30 (868), B11 (822), B91 (583), A15 (780), B26 (808), C47 (783). Let me start by selecting these. Now, let's calculate their total effort:} \\
\\ 
\texttt{A30: 515} \\
\texttt{B11: 571} \\
\texttt{B91: 252} \\
\texttt{A15: 803} \\
\texttt{B26: 493} \\
\texttt{C47: 211} \\
\texttt{Adding these up: 515 + 571 = 1086; 1086 + 252 = 1338; +803 = 2141; +493 = 2634; +211 = 2845. That's 2845. Under 4379, so there's room for more.} \\
\\ 
\texttt{Next, check the next highest importance items. A75 (461) has effort 401. Adding that: 2845 + 401 = 3246. Then A39 (484) with 980? Wait, 980 is a high effort. 3246 + 980 = 4226. Still under 4} \\
            }
        }
    }
    & & \\ \\

    \theutterance \stepcounter{utterance}  
    & & \multicolumn{4}{p{0.6\linewidth}}{
        \cellcolor[rgb]{0.9,0.9,0.9}{
            \makecell[{{p{\linewidth}}}]{
                \texttt{\tiny{[P1$\langle$GM]}}
                \texttt{Your response did not contain an argument. You must include the structured format ARGUMENT: \{...\} somewhere in your response. Try again} \\
            }
        }
    }
    & \\ \\

    \theutterance \stepcounter{utterance}  
    & \multicolumn{4}{p{0.6\linewidth}}{
        \cellcolor[rgb]{0.9,0.9,0.9}{
            \makecell[{{p{\linewidth}}}]{
                \texttt{\tiny{[P1$\rangle$GM]}}
                \texttt{<think>} \\
\texttt{Okay, let me try again. I need to make sure I include an argument in my response. So, first, I'll outline my initial proposal as before, selecting the highest importance items. Then, I need to structure an argument that explains why this proposal is beneficial. The argument should mention the total effort and how it's under the limit, and maybe highlight the high importance items. Let me check the numbers again to confirm. The total effort was 4226, which is under 4379. The importance sum includes A30, B11, B91, A15, B26, C47, A75, and A39. Wait, A39's effort is 980, so adding that to 3246 (previous total) gives 4226. Yes. The argument should state that this set maximizes my importance while staying within the limit. Also, maybe mention that removing any of these would lower the total importance, so it's optimal. Need to phrase it in a way that's persuasive but doesn't reveal my exact importance values. Maybe say "high{-}value items" instead of specific numbers. Alright, let's put that into the argument.} \\
\texttt{</think>} \\
\\ 
\texttt{PROPOSAL: \{"A30", "B11", "B91", "A15", "B26", "C47", "A75", "A39"\}  } \\
\texttt{ARGUMENT: \{"The proposed set maximizes my priority items while staying well under the effort limit (total effort: 4226). Removing any item would significantly reduce my total importance, as these represent the highest{-}value items in my evaluation."\}} \\
            }
        }
    }
    & & \\ \\

    \theutterance \stepcounter{utterance}  
    & & \multicolumn{4}{p{0.6\linewidth}}{
        \cellcolor[rgb]{0.9,0.9,0.9}{
            \makecell[{{p{\linewidth}}}]{
                \texttt{\tiny{[P1$\langle$GM]}}
                \texttt{Your response contained an untagged sequence or you used STRATEGIC REASONING more than once. You may only use the structured formats as explained in the initial message. They must all be of the form TAG: \{...\}} \\
            }
        }
    }
    & \\ \\

    \theutterance \stepcounter{utterance}  
    & \multicolumn{4}{p{0.6\linewidth}}{
        \cellcolor[rgb]{0.9,0.9,0.9}{
            \makecell[{{p{\linewidth}}}]{
                \texttt{\tiny{[P1$\rangle$GM]}}
                \texttt{<think>} \\
\texttt{Okay, let me try to fix this. The user said my previous response had an untagged sequence or used strategic reasoning more than once. I need to make sure all my messages strictly follow the structured formats: PROPOSAL, REFUSE, ARGUMENT, or AGREE. Each must be properly tagged with a colon and whitespace, and the content should be a set or a string in the case of ARGUMENT.} \\
\\ 
\texttt{In my last attempt, I wrote PROPOSAL with a set and then ARGUMENT with a string. The user's error message suggests that maybe I included some unstructured text. Wait, looking back, I see that after the ARGUMENT tag, I used curly braces with a string. But the ARGUMENT should be a set containing the argument as a single element? Or is the argument supposed to be a string within the set? Wait, the initial instructions say: "The argument must be a python set containing 0 or more strings." So the ARGUMENT should be a set of strings. So in my previous response, I had ARGUMENT: \{"The proposed set..."\}, which is a set with one string. That should be correct. But maybe the system is expecting the argument to be a single string without the quotes? Or perhaps the formatting is off.} \\
\\ 
\texttt{Wait, the user's example shows that the argument is a set containing the string. So the correct format is ARGUMENT: \{"This is my argument."\}. So my previous response was correct in that aspect. Hmm. The user's error message says "your response contained an untagged sequence or you used STRATEGIC REASONING more than once." Maybe I included some text outside the structured formats. Let me check my last message again.} \\
\\ 
\texttt{Original response:} \\
\texttt{PROPOSAL: \{"A30", "B11", "B91", "A15", "B26", "C47", "A75", "A39"\}  } \\
\texttt{ARGUMENT: \{"The proposed set maximizes my priority items while staying well under the effort limit (total effort: 4226). Removing any item would significantly reduce my total importance, as these represent the highest{-}value items in my evaluation."\}} \\
\\ 
\texttt{Wait, the PROPOSAL line has a space after the colon, which is correct. The ARGUMENT line also has a space after the colon. The content is a set with a single string. So maybe the issue is that the system is expecting} \\
            }
        }
    }
    & & \\ \\

    \theutterance \stepcounter{utterance}  
    & & & \multicolumn{2}{p{0.3\linewidth}}{
        \cellcolor[rgb]{0.9,0.9,0.9}{
            \makecell[{{p{\linewidth}}}]{
                \texttt{\tiny{[GM$|$GM]}}
                \texttt{abort} \\
            }
        }
    }
    & & \\ \\

    \theutterance \stepcounter{utterance}  
    & & & \multicolumn{2}{p{0.3\linewidth}}{
        \cellcolor[rgb]{0.9,0.9,0.9}{
            \makecell[{{p{\linewidth}}}]{
                \texttt{\tiny{[GM$|$GM]}}
                \texttt{end game} \\
            }
        }
    }
    & & \\ \\

\end{supertabular}
}

\end{document}
